
%\needspace{4\baselineskip}
\color{unidarkgreen}{\subsection{Фиксация положения выключателей В3 и ШСВ/СВ (ФПВ В3, ФПВ ШСВ/СВ)}}
\color{black}

\begin{enumerate}[label=\arabic{section}.\arabic{subsection}.\arabic{enumi}, labelsep=4pt, leftmargin=0pt, itemindent=57pt, itemsep=0pt, parsep=0pt] % если что можно крутить itemsep

\item
Функциональный блок фиксации положения выключателей (ФПВ) предназначен для определения состояния выключателей на основании технологической информации:

\begin{itemize}
    \item состояния блок-контактов выключателя;
    \item положения разъединителей в цепи выключателя;
    \item информации о выводе в ремонт выключателя.
\end{itemize}

\item
Контроль положения разъединителей при формировании сигналов включенного или отключенного положения выключателей (например, \textit{«В3 включен»} и \textit{«ШСВ/СВ отключен»} для выключателя ШСВ/СВ) вводится программным ключом \textbf{«Контр. Р В3»} (\textbf{«Контр. Р ШСВ/СВ»} для выключателя ШСВ/СВ).

\item
Функция ФПВ пофазно контролирует состояние блок-контактов выключателя и в случае расхождения полюсов формирует сигнал \textit{«Внутр. пуск. ЗНФ»}.

\item
Функциональная схема логики ФПВ В3 приведена на рисунке \ref{pict_fpv:fpv3}. Уставки ФПВ В3 представлены в таблице \ref{app_set:fpv3}.

\vspace{2mm}
\begin{figure}[H]
\centering
\includegraphics[width=1\textwidth,height=1\textheight,keepaspectratio]{\fbpath/761.3461 ФПВ34/img/fpv3.pdf}
\captionof{figure}{Функциональная схема логики ФПВ В3}\label{pict_fpv:fpv3}
\end{figure}


\item
Функциональная схема логики ФПВ ШСВ/СВ приведена на рисунке \ref{pict_fpv:fpv4}. Уставки ФПВ ШСВ/СВ представлены в таблице \ref{app_set:fpv4}.

\vspace{2mm}
\begin{figure}[H]
\centering
\includegraphics[width=1\textwidth,height=1\textheight,keepaspectratio]{\fbpath/761.3461 ФПВ34/img/fpv4.pdf}
\captionof{figure}{Функциональная схема логики ФПВ ШСВ/СВ}\label{pict_fpv:fpv4}
\end{figure}



\end{enumerate}